\section*{Finite-support compression for fixed typed projected trees}
\label{sec:fixed-tree-support-compression}

Consider a fixed finite typed tree $T$ in the repository's convention: each
type $\tau$ has one finite-dimensional Hilbert space $V_\tau$ and one
orthogonal projector $P_\tau$, each leaf supplies a fixed vector, and each
internal node has a bounded multilinear law.  Let $\ell_\tau$ be the number
of leaves of type $\tau$ and $k_\tau$ the number of internal nodes whose
output type is $\tau$.

\textbf{Theorem (support compression).}  Every admissible finite-dimensional
configuration has an evaluation-preserving compressed configuration with
\[
  \dim W_\tau\le B_\tau:=2(\ell_\tau+2k_\tau),
  \qquad \operatorname{rank}P_{W_\tau}\le B_\tau,
\]
for every type $\tau$.  The compressed laws have operator norms no larger
than the original laws, and their projected-input closure residual norms are
no larger than the original residual norms.  All ambient, recursively
projected, reduced, normal, and projected-root errors are exactly preserved.

\textbf{Construction.}  For every type $\tau$, collect the leaf vectors of
that type and, for every node outputting in that type, collect two vectors:
the ambient node value
\[
  F_v=\mu_v(F_{v_1},\ldots,F_{v_a})
\]
and the unprojected value in the recursively projected evaluation
\[
  H_v=\mu_v(R_{v_1},\ldots,R_{v_a}),
  \qquad R_v=P_\tau H_v.
\]
Let $W_\tau$ be the smallest subspace containing these vectors and invariant
under $P_\tau$.  The invariant span is generated by the original vectors and
their $P_\tau$ images, so
\[
  \dim W_\tau\le2(\ell_\tau+2k_\tau).
\]

Let $\Pi_\tau$ be the orthogonal projector onto $W_\tau$.  Define
\[
  P'_\tau=P_\tau|_{W_\tau},
  \qquad
  \mu'_v=\Pi_{\tau_{out}}\,\mu_v|_{W_{\tau_1}\times\cdots\times W_{\tau_a}}.
\]
Because $W_\tau$ is $P_\tau$-invariant, $\Pi_\tau$ commutes with
$P_\tau$.  Induction over the tree shows that the compressed ambient values
are exactly the original $F_v$ and the compressed raw projected values are
exactly the original $H_v$; hence every projected value and every named root
error is unchanged.

For the norm bounds, restriction and orthogonal output compression are
contractions, so $\|\mu'_v\|_{op}\le\|\mu_v\|_{op}$.  Moreover,
\[
 (I-P'_\tau)\mu'_v(P'\cdot,\ldots,P'\cdot)
 =\Pi_\tau(I-P_\tau)\mu_v(P\cdot,\ldots,P\cdot)
\]
on the compressed input spaces.  Its norm is therefore no larger than the
original closure residual norm.  If the admissible convention requires a
proper nonzero projector, append one unused range coordinate and one unused
kernel coordinate and extend all laws by zero; this increases the dimension
bound by at most two without changing any evaluation or norm cap.

\textbf{Corollary (fixed-tree dimension/rank reduction).}  The supremum over
all finite ambient dimensions and projector ranks for a fixed finite typed
tree equals the supremum over the bounded dimensions above (or the padded
bound when proper projectors are required).  For a one-type binary $k=3$
chain or branching tree, $\ell=4$, $k=3$, so the bound is $20$ (or $22$ with
proper-projector padding).

\textbf{Corollary (global $k=3$ strict supremum gap).}  Fix $0<\eta<1$ and
the independent-law binary $k=3$ chain or branching class.  The compression
corollary reduces its global supremum to a finite union of compact parameter
spaces: bounded tensors, unit leaves, and fixed-rank orthogonal projectors.
The projected error is continuous, so the supremum is attained.  M10 proves
that no admissible configuration attains $U_3(\eta)$; consequently
\[
  C_{3,\mathrm{ind}}^P(\eta)<U_3(\eta),\qquad 0<\eta<1.
\]
This is a strict but non-explicit gap; determining its exact value remains a
separate open problem.  The theorem uses the repository's typed-projector
convention and does not assert a single dimension bound uniform over trees
whose size or arity is allowed to grow without bound.
